\section{Limitations, trade-offs and production changes}

\subsection{Limitations}
\begin{itemize}
\item \textbf{Single nodes.} Kafka has one broker with replication factor 1, and Spark runs in local mode. A
broker failure would stop ingestion, and one machine limits the throughput.
\item \textbf{Simple model.} Vehicles move in straight lines, and zones are nearest-centre areas on a coarse
grid, so a point near a zone border can fall in the wrong zone.
\item \textbf{Estimated lateness.} The batch layer counts an event as late when it arrived more than
10 simulated minutes after it happened. This is close to, but not exactly, what the speed layer dropped.
\item \textbf{Light tracing and local security.} We trace with ids and lineage tables instead of OpenTelemetry
and Jaeger. Passwords are in an \code{.env} file, and Kafka and PostgreSQL use no TLS.
\item \textbf{Limited memory.} The full stack needs most of the memory that Docker can use on an 8\,GB
laptop, so heavy desktop apps make the batch runs slower.
\end{itemize}

\subsection{Trade-offs we accepted}
\begin{itemize}
\item \textbf{Two code paths} give exact money numbers and cheap recomputing. Shared code and a parity test
reduce the risk that they drift apart.
\item \textbf{A 10-second trigger} for the aggregate queries instead of 5 seconds halves the CPU they use.
\item \textbf{A 10-minute watermark} keeps state small and dashboards fresh, but the speed layer drops some
late events. The batch layer counts them after midnight.
\item \textbf{Collecting to the driver.} The \code{vehicle\_state} query collects each small micro-batch
(about 100 rows) to the Spark driver. This is much faster than several Spark jobs per batch, and the input
cap of 5000 offsets per trigger keeps it safe.
\item \textbf{Kafka native image.} It saves memory, but Apache says it is for development and testing. The
Java image is a drop-in replacement.
\end{itemize}

\subsection{What we would change at production scale}
\begin{itemize}
\item \textbf{Kafka:} a cluster of at least three brokers with replication factor 3, and a schema registry
with Avro or Protobuf schemas for the events.
\item \textbf{Spark:} a cluster (Kubernetes or a managed service) with more partitions and the RocksDB state
store for large state.
\item \textbf{Master dataset:} S3 or HDFS with an open table format such as Delta Lake or Apache Iceberg,
and regular compaction of the small files that streaming creates.
\item \textbf{Airflow:} the Kubernetes or Celery executor, with Spark batch jobs started on the cluster.
\item \textbf{Zones:} a real geo index, for example H3 cells or zone polygons.
\item \textbf{Operations:} central logs (Loki or Elasticsearch), OpenTelemetry tracing, alerts by e-mail,
SMS or Slack, a secret manager and TLS everywhere.
\end{itemize}
